% !TEX program = xelatex
% ===========================================================================
%  《基本不等式》—— 中文数学习题课讲义（Loom 模板 · 课前自学版）
%  必须用 XeLaTeX 编译两遍；QA 用 scripts/qa_check.py。
%  第一行的魔法注释让 VS Code (LaTeX Workshop) 走 XeLaTeX——别删。
% ===========================================================================
\documentclass{loom}

% 挑战题 ★ 用 Zapf Dingbats 绘制，不依赖系统字体是否含 U+2605
\usepackage{pifont}

% 页眉线（可选）；偏好极简时可整行注释掉
%\runningthread{《基本不等式》 · 课前自学讲义}

% ---------------------------------------------------------------------------
%  讲义专用组件：题面（茜红结）、解析（靛蓝结）、四段式解答标签
% ---------------------------------------------------------------------------
% 题目框：breakable=false 保证一题不跨页。默认不留答题空白区。
% —— 若要做「打印作答版」，取消下面 \answerspace 定义与 problem 框 after upper
%    一行的注释即可，题末会追加一块纯白书写面板（默认 5 行，可 \answerspace[8]）。
% \newcommand{\answerspace}[1][5]{\par\medskip\noindent
%   \colorbox{white}{\parbox[t][#1\baselineskip][t]{\dimexpr\linewidth-2\fboxsep\relax}{\strut}}}
\newtcolorbox{problem}[1]{knotmadder, breakable=false,
  before upper={{\headingfont\bfseries\color{madder}#1}\par\smallskip}%
  % , after upper={\answerspace}   % 打印作答版：取消本行注释
}
\newtcolorbox{solution}[1]{knotindigo,
  before upper={{\headingfont\bfseries\color{indigo}#1}\par\smallskip}}

\newcommand{\think}{\par\smallskip\noindent
  {\headingfont\bfseries\color{indigo}〔怎么想〕}\enspace}
\newcommand{\solve}{\par\smallskip\noindent
  {\headingfont\bfseries\color{inkiron}〔解〕}\enspace}
\newcommand{\solvehead}[1]{\par\smallskip\noindent
  {\headingfont\bfseries\color{inkiron}〔#1〕}\enspace}
\newcommand{\checkup}{\par\smallskip\noindent
  {\headingfont\bfseries\color{selvage}〔检验〕}\enspace}
\newcommand{\insight}[1]{\par\smallskip\noindent
  {\headingfont\bfseries\color{weld!60!inkiron}〔方法小结〕}\enspace
  {\strandfont\small #1}\par}
\newcommand{\hint}[1]{\par\smallskip\noindent
  {\strandfont\small{\raisebox{-0.8pt}{\loomtile{weld}}\hspace{4pt}}%
  {\bfseries\color{weld!60!inkiron}入手提示}\enspace{\color{inkiron!85!white}#1}}\par}
% 页边旁批
\newcommand{\sidenote}[1]{\marginnote{\raggedright\strandfont\footnotesize
  \color{selvage}#1}}
\newcommand{\starred}{{\color{weld!70!inkiron}\ding{72}}}
% 记号 (☆)：a^2+b^2 >= 2ab 的全文代号
% 空心星用 pifont \ding{73} 绘制——字面 ☆(U+2606) 在部分拉丁字体里没有字形。
\newcommand{\starref}{(\text{\ding{73}})}

\setlist[enumerate]{leftmargin=2.3em, itemsep=2pt, topsep=3pt}
\setlist[itemize]{leftmargin=1.6em, itemsep=2pt, topsep=3pt}

\begin{document}

\loomcover
  {基本不等式}
  {}
  {}
  {二〇二六年七月}

\clearpage
{\small\tableofcontents}
\clearpage

% ---- 第零章：预备知识 ----
% ===========================================================================
%  第零章 · 预备知识
% ===========================================================================
\setcounter{section}{-1}
\section{预备知识}

本章列出正文将反复使用的三条预备知识。

\block{0.1 比较大小的基本方法：作差}

比较实数 $a$、$b$ 的大小，最基本的方法是考察差 $a-b$ 的符号：
\[
a-b>0 \iff a>b,\qquad a-b=0 \iff a=b,\qquad a-b<0 \iff a<b.
\]
全文所有比较大小、证明不等式的论证，最终都归结为这一条。作差法的长处是便于通分，下面是一个典型例子。

\begin{proposition}[糖水不等式]
用 $a$ kg 白糖制出 $b$ kg 糖水（$b>a>0$），浓度为 $\dfrac ab$；再加入 $m$ kg 白糖（$m>0$），浓度变为 $\dfrac{a+m}{b+m}$。经验告诉我们糖水加糖后变甜，即
\[
\frac{a+m}{b+m}>\frac ab.
\]
\end{proposition}

\begin{proof}
作差、通分：
\[
\frac{a+m}{b+m}-\frac ab=\frac{b(a+m)-a(b+m)}{b(b+m)}=\frac{m(b-a)}{b(b+m)},
\]
分子、分母的每个因子均为正数，故差为正，不等式成立。（书例 1.1-2）
\end{proof}

\sidenote{传递性：由 $A\ge B$ 与 $B\ge C$ 可得 $A\ge C$。本讲中多次以「链式放缩」的形式出现（题 3）。}

\block{0.2 平方的非负性}

由完全平方公式 $(a\pm b)^2=a^2\pm 2ab+b^2$，立即得到本讲义最基本的事实：
\[
(a-b)^2\ge 0,\quad\text{当且仅当 } a=b \text{ 时取等号}.
\]
全文的不等式都以此为源头。另记一个后面要用的恒等式：
\[
(a-b)^2+4ab=(a+b)^2.
\]

\block{0.3 圆周角定理与射影定理}

以下两条定理在板块二与拓展 A 中直接使用，证明要点附后。

\begin{proposition}[圆周角定理（直径情形）]
设 $AB$ 是圆 $O$ 的直径，$D$ 为圆上异于 $A$、$B$ 的点，则 $\angle ADB=90^\circ$。
\end{proposition}

\begin{proof}[证明要点]
$OA=OD=OB$，故 $\triangle OAD$、$\triangle OBD$ 均为等腰三角形，$\angle ODA=\angle OAD$，$\angle ODB=\angle OBD$。在 $\triangle ABD$ 中，
$\angle ADB=\angle ODA+\angle ODB=\angle OAD+\angle OBD$，代入内角和 $180^\circ$，得 $2\angle ADB=180^\circ$，即 $\angle ADB=90^\circ$。
\end{proof}

\begin{proposition}[射影定理（高的情形）]
直角三角形斜边上的高，是垂足所分斜边两段的比例中项：设高为 $h$，两段长为 $p$、$q$，则
\[
h^2=pq.
\]
\end{proposition}

\begin{proof}[证明要点]
高把原三角形分成两个小直角三角形，由锐角互余知两个小三角形相似，对应边成比例 $\dfrac hp=\dfrac qh$，即得 $h^2=pq$。
\end{proof}

\clearpage


% ---- 第一部分：题目区 ----
% ===========================================================================
%  板块一 · 基本不等式及其变形（卡片 1 + 题 1–4）
% ===========================================================================
\section{基本不等式及其变形}

\begin{strand}
\keyword{本板块主线}：由 \starref{} 经替换得到基本不等式，认清「和」与「积」两种形态，并明确取等条件。
\end{strand}

\block{引例：连续打折的平均折扣}

一件衣服先打 8 折，再打 9 折。能说「平均每次打 8.5 折」吗？

两次连打后价格变为原来的 $0.8\times 0.9=0.72$。「平均每次打 $x$ 折」应当满足 $x\cdot x=0.72$，即
\[
x=\sqrt{0.72}\approx 0.849,
\]
不是 $0.85$。注意这里的「平均」是对\textbf{乘法}说的：把两次不同的折扣换成两次相同的折扣，总效果不变。一般地：
\begin{itemize}
  \item 把 $a+b$ 换成两个相同数之和：$a+b=\dfrac{a+b}{2}+\dfrac{a+b}{2}$，数 $\dfrac{a+b}{2}$ 称为 $a$、$b$ 的\keyword{算术平均数}；
  \item 把 $ab$ 换成两个相同数之积：$ab=\sqrt{ab}\cdot\sqrt{ab}$，数 $\sqrt{ab}$ 称为正数 $a$、$b$ 的\keyword{几何平均数}。
\end{itemize}
打折问题里 $\sqrt{0.72}\approx 0.849<0.85=\dfrac{0.8+0.9}{2}$：几何平均数小于算术平均数。这不是巧合，下面给出一般结论。

\block{知识精讲：从 \starref{} 到基本不等式}

对任意实数 $a$、$b$，把 $(a-b)^2\ge 0$ 展开移项，得到本讲的出发点，记作 \starref{}：
\[
a^2+b^2\ge 2ab\quad(a,b\in\mathbb{R}),\ \text{当且仅当 } a=b \text{ 时取等号}. \tag{\ding{73}}
\]

设 $a>0$，$b>0$。在 \starref{} 中用 $\sqrt a$ 替换 $a$、用 $\sqrt b$ 替换 $b$（替换要求 $\sqrt a$、$\sqrt b$ 有意义，这正是限定 $a,b>0$ 的原因）：
\[
(\sqrt a)^2+(\sqrt b)^2\ge 2\sqrt a\cdot\sqrt b,\qquad\text{即}\qquad a+b\ge 2\sqrt{ab},
\]
两边除以 2，得

\begin{proposition}[基本不等式]
\[
\frac{a+b}{2}\ \ge\ \sqrt{ab}\quad(a,b>0),\ \text{当且仅当 } a=b \text{ 时取等号}.
\]
\end{proposition}

取等条件由 \starref{} 继承而来：$\sqrt a=\sqrt b$，即 $a=b$。这意味着：\textbf{两个正数的算术平均数，不小于它们的几何平均数。}打折问题不是巧合，是定理。%
\sidenote{另一条推导：由 0.2 的恒等式 $(a-b)^2+4ab=(a+b)^2$ 得 $(a+b)^2\ge 4ab$；$a,b>0$ 时两边均非负，开方即得 $a+b\ge 2\sqrt{ab}$。（书 1.2.1）}

\textbf{使用步骤（比大小、放缩时）。}
\begin{enumerate}
  \item 确认两个量\textbf{都是正数}；
  \item 认出「和」与「积」两种形态：$a+b\ge 2\sqrt{ab}$，或 $ab\le\left(\dfrac{a+b}{2}\right)^2$；
  \item 写清\textbf{等号何时成立}（把 $a=b$ 代回检验）。
\end{enumerate}

\textbf{常用变形一览}（$a,b>0$）：
% 注释开关：自总结版可把下表整表留空，让学生课后自己填
\begin{center}
\begin{tabularx}{\linewidth}{@{}l l L@{}}
\toprule
\multicolumn{1}{l}{\color{indigo}变形} &
\multicolumn{1}{l}{\color{indigo}形式} &
\multicolumn{1}{l}{\color{indigo}记忆} \\
\midrule
和的下界 & $a+b\ge 2\sqrt{ab}$                  & 和 $\ge$ 两倍根号积 \\
积的上界 & $ab\le\left(\dfrac{a+b}{2}\right)^2$ & 积 $\le$ 半和的平方 \\
倒数对   & $\dfrac{b}{a}+\dfrac{a}{b}\ge 2$     & 互倒之和不小于 2 \\
特例     & $a+\dfrac{1}{a}\ge 2$                & $x$ 与 $\dfrac1x$（$x>0$） \\
\bottomrule
\end{tabularx}
\end{center}

% 注释开关：易错提醒（自总结版可 \iffalse … \fi 注释）
\textbf{易错提醒。}
\begin{itemize}
  \item 基本不等式的\textbf{前提是正数}。$a+\dfrac1a\ge 2$ 对 $a<0$ 不成立（此时 $a+\dfrac1a\le -2$，见下面的「想一想」）。
  \item \starref{} 对\textbf{一切实数}成立；基本不等式只对\textbf{正数}成立，二者适用范围不同。特别地，$ab\le\left(\dfrac{a+b}{2}\right)^2$ 对一切实数成立（它就是 \starref{} 的移项配方），这一点在拓展 D 中是关键。
\end{itemize}

\begin{remark*}
想一想：\textbf{(1)} $\left(\dfrac{a+b}{2}\right)^2\ge ab$ 与 $\dfrac{a+b}{2}\ge\sqrt{ab}$ 等价吗？（前者对一切实数成立；后者含根号，只在 $a,b\ge 0$ 时有意义。）\quad
\textbf{(2)} $a+\dfrac{1}{a}\ge 2$（$a\ne 0$）恒成立吗？取 $a=-1$ 检验。
\end{remark*}

\begin{problem}{题 1 （作差与公式的选择）}
已知 $a,b>0$，比较下列各组数的大小，并说明理由：

(1) $\dfrac{a}{b}+\dfrac{b}{a}$ 与 $2$；\quad
(2) $a^2+3$ 与 $3a$（本小问 $a$ 为任意实数）；\quad
(3) 把 (2) 中的 $3$ 都换成 $2$，结论有何变化？
\hint{(1) 是「互倒之和」；(2) $3a$ 可能为负，基本不等式的正数前提不满足时，回到 0.1 的作差法，对差配方。}
\end{problem}

\begin{problem}{题 2 （四数排序）}
设 $0<a<b$，把 $a,\ b,\ \sqrt{ab},\ \dfrac{a+b}{2}$ 从小到大排列，并逐段证明。
\hint{三个空隙，三次比较：$a$ 与 $\sqrt{ab}$ 比较时提出公因式 $\sqrt{a}$；中间一段是基本不等式（$a\ne b$，等号取不到，写严格不等号）。}
\end{problem}

\begin{problem}{题 3 \starred（两步放缩）}
已知 $a>0$，比较 $a^2+\dfrac{2}{a}$ 与 $3$ 的大小。
\hint{一次基本不等式不够用（一项带平方、一项带分母）。分两步：先用 $a^2+1\ge 2a$ 降次，再处理剩下的部分；用 0.1 的传递性衔接，并检查两步的等号能否\textbf{同时}成立。}
\end{problem}

\begin{problem}{题 4 （下界与上界）}
已知 $m=a+\dfrac{1}{a-2}$（$a>2$），$n=4-b^2$（$b\ne 0$），比较 $m$ 与 $n$ 的大小。
\hint{把 $a$ 拆成 $(a-2)+2$，$m$ 中出现互为倒数的一对；再考察 $n$ 的上界是多少、取得到吗。}
\end{problem}

\clearpage

% ===========================================================================
%  板块二 · 基本不等式的几何解释（卡片 2 + 题 5）
% ===========================================================================
\section{基本不等式的几何解释}

\begin{strand}
\keyword{本板块主线}：半弦不过半径。
\end{strand}

\block{知识精讲：构图}

\starref{} 有一张著名的图：\keyword{赵爽弦图}。四个直角三角形（直角边 $a$、$b$）拼入大正方形，四个三角形面积之和 $2ab$ 不超过大正方形面积 $a^2+b^2$，空隙恰是 $(a-b)^2$。

基本不等式 $\dfrac{a+b}{2}\ge\sqrt{ab}$ 说的则是\textbf{长度}。构图按以下问题展开（题 5 完整完成这条路）：

\begin{itemize}
  \item 问题 1：$\dfrac{a+b}{2}$ 与 $\sqrt{ab}$ 都是长度——把长为 $a$、$b$ 的两条线段首尾相接：在直线上取 $A$、$E$、$B$，使 $AE=a$，$EB=b$；取 $AB$ 的中点 $O$，则 $OA=OB=\dfrac{a+b}{2}$。
  \item 问题 2：长为 $\dfrac{a+b}{2}$ 的线段有多少条？——以 $O$ 为圆心、$\dfrac{a+b}{2}$ 为半径作圆，每条半径都是。
  \item 问题 3：$\sqrt{ab}$ 在哪里？——过 $E$ 作 $AB$ 的垂线，交半圆于 $D$，连接 $AD$、$BD$。由圆周角定理（0.3），$\angle ADB=90^\circ$；由射影定理（0.3），$DE^2=AE\cdot EB=ab$，即 $DE=\sqrt{ab}$。
  \item 结论：$DE$ 是过 $E$ 的半弦，$DO$ 是半径，半弦不超过半径，即 $\sqrt{ab}\le\dfrac{a+b}{2}$。
\end{itemize}

\begin{center}
\begin{tikzpicture}[scale=1.05, line cap=round]
  % 半圆：直径 AB，AE=a=1.4, EB=b=3.6，半径 r=2.5
  \coordinate (A) at (0,0);
  \coordinate (B) at (5,0);
  \coordinate (O) at (2.5,0);
  \coordinate (E) at (1.4,0);
  \coordinate (D) at (1.4,2.2449); % sqrt(1.4*3.6)=sqrt(5.04)
  \draw[inkiron] (A) -- (B);
  \draw[inkiron] (B) arc (0:180:2.5);
  \draw[selvage] (A) -- (D) -- (B);
  \draw[madder, very thick] (E) -- (D) node[midway, left, madder]{\small $\sqrt{ab}$};
  \draw[indigo, very thick] (O) -- (D) node[midway, above right=-1pt, indigo]{\small $\dfrac{a+b}{2}$};
  % 直角记号
  \draw[inkiron] (1.4,0.16) -- (1.56,0.16) -- (1.56,0);
  % 点与标注
  \foreach \p/\n/\where in {A/A/below, B/B/below, O/O/below, E/E/below, D/D/above}{
    \fill[inkiron] (\p) circle (1.3pt);
    \node[\where, inkiron] at (\p) {\small $\n$};
  }
  \draw[weld, |-|] (0,-0.55) -- node[below, weld]{\small $a$} (1.4,-0.55);
  \draw[weld, |-|] (1.4,-0.55) -- node[below, weld]{\small $b$} (5,-0.55);
\end{tikzpicture}
\end{center}

\textbf{等号的位置。}半弦等于半径，只能是垂足 $E$ 与圆心 $O$ 重合，此时 $a=b$，$D$ 位于半圆最高点，$DE$ 本身就是一条半径。「当且仅当 $a=b$ 取等号」在图上对应一个确定的位置。

\begin{problem}{题 5 （半圆中的构造）}
在直线上依次取点 $A$、$E$、$B$，使 $AE=a$，$EB=b$（$a,b>0$）。以 $AB$ 为直径、中点 $O$ 为圆心作半圆，过 $E$ 作 $AB$ 的垂线交半圆于 $D$，连接 $AD$、$BD$、$DO$。

(1) 说明 $\angle ADB=90^\circ$，并用射影定理求 $DE$（用 $a$、$b$ 表示）；

(2) 在 $\mathrm{Rt}\triangle DEO$ 中比较 $DE$ 与 $DO$，写出对应的不等式；

(3) 等号成立时，$E$、$D$ 各在什么位置？此时 $a$、$b$ 满足什么关系？
\hint{(1) 0.3 的两条定理各用一次；(2) 斜边与直角边孰长；(3) 考察 $DE=DO$ 时三角形退化成什么。}
\end{problem}

\clearpage

% ===========================================================================
%  板块三 · 用基本不等式证明不等式（卡片 3 + 题 6–7）
% ===========================================================================
\section{用基本不等式证明不等式}

\begin{strand}
\keyword{本板块主线}：把待证不等式拆成若干基本不等式，逐项放缩；各项取等条件必须能同时成立。
\end{strand}

\block{知识精讲：逐项放缩}

证明「整体 $\ge$ 整体」的常用步骤：
\begin{enumerate}
  \item \textbf{拆}：把左边拆成若干个可以使用「和 $\ge$ 两倍根号积」的部分；
  \item \textbf{乘或加}：各部分的不等式同向相乘（\textbf{均为正数才可相乘}）或相加；
  \item \textbf{对齐取等条件}：每一部分取等号的条件必须能同时成立，否则合成后的等号取不到，只能写严格不等号。本讲把这一步称为「\keyword{会师}」。
\end{enumerate}

模板例：$x,y>0$，证 $(x+y)\left(\dfrac1x+\dfrac1y\right)\ge 4$。

拆：$x+y\ge 2\sqrt{xy}$，$\dfrac1x+\dfrac1y\ge 2\sqrt{\dfrac{1}{xy}}$；两式两边均为正，相乘得左边 $\ge 4$。取等条件都是 $x=y$，可同时成立。%
\sidenote{若改证 $(x+y)\left(\frac1x+\frac4y\right)\ge 9$，两部分的取等条件一个是 $x=y$、一个是 $y=2x$，无法同时成立，此路不通；应先展开、再对剩余项用基本不等式。见课后练习 4。}

\begin{problem}{题 6 （逐项放缩相乘）}
已知 $x,y>0$，求证：
\[
(x+y)(x^2+y^2)(x^3+y^3)\ge 8x^3y^3.
\]
\hint{右边 $8=2\times2\times2$，提示左边三个括号各用一次基本不等式。三处等号条件是否一致？}
\end{problem}

\begin{problem}{题 7 \starred（用条件替换 1）}
已知 $a,b,c>0$，且 $a+b+c=1$，求证：
\[
\frac{1}{a}+\frac{1}{b}+\frac{1}{c}\ge 9.
\]
\hint{把分子上的 $1$ 全部替换为 $a+b+c$，展开后清点：3 个常数 1，加上几对「互倒之和」？这一替换手法是板块五的核心。}
\end{problem}

\clearpage

% ===========================================================================
%  板块四 · 求最值：一正二定三相等（卡片 4 + 题 8–13）
% ===========================================================================
\section{求最值：一正二定三相等}

\begin{strand}
\keyword{本板块主线}：用不等式求最值，成立与否取决于等号能否取到。
\end{strand}

\block{知识精讲：不等式为什么能求最值}

此前求最值有两条路：几何方法（如「将军饮马」，三点共线时最短）与函数方法（如二次函数配方）。基本不等式提供第三条路，其逻辑必须想清楚：

\begin{strand}
$y\ge 4$ 只说明 $y$ 不会小于 4，\textbf{不等于}说「$y$ 的最小值是 4」——$y=5$ 也满足 $y\ge 4$。只有当 $y=4$ 确实取得到（等号成立），4 才是最小值。
\end{strand}

因此用基本不等式求最值要做三项检查，\keyword{口诀：一正、二定、三相等}：

\begin{center}
\begin{tabularx}{\linewidth}{@{}l L@{}}
\toprule
\multicolumn{1}{l}{\color{indigo}检查} &
\multicolumn{1}{l}{\color{indigo}原因} \\
\midrule
\textbf{一正}：各项为正 & 基本不等式的前提 \\
\textbf{二定}：和或积为定值 & $x+y\ge 2\sqrt{xy}$ 的右边必须是常数才构成下界；积 $xy=P$ 为定值 $\Rightarrow$ 和的最小值 $2\sqrt{P}$；和 $x+y=S$ 为定值 $\Rightarrow$ 积的最大值 $\dfrac{S^2}{4}$ \\
\textbf{三相等}：等号取得到 & 取不到则该界不是最值（题 13 专门辨析） \\
\bottomrule
\end{tabularx}
\end{center}

与板块二的图对照：半弦 $DE=\sqrt{ab}$ 何时最长？垂足 $E$ 与圆心重合、$a=b$ 之时。\textbf{最值出现在等号成立的位置}，几何与代数一致。

% 注释开关：易错提醒（自总结版可 \iffalse … \fi 注释）
\textbf{易错提醒。}
\begin{itemize}
  \item 右边不是定值时，「$\ge$」不构成最小值结论。$x^2+\dfrac1x\ge 2\sqrt{x}$ 成立，但 $2\sqrt{x}$ 随 $x$ 变化，从中读不出最小值（题 13 D 即此错误）。
  \item 负变量先转正：$x<0$ 时对 $-x>0$ 使用基本不等式（题 9）。
  \item 取等的 $x$ 必须落在定义域内，代回检验是必要步骤。
\end{itemize}

\begin{problem}{题 8 （直接使用）}
当 $x>0$ 时，求 $y=4x+\dfrac{12}{x}$ 的最小值，并指出取到最小值时的 $x$。
\hint{两项之积是否为定值？计算 $4x\cdot\dfrac{12}{x}$。}
\end{problem}

\begin{problem}{题 9 （负变量先转正）}
已知 $x<0$，求 $y=x+\dfrac{1}{2x}-2$ 的最大值。
\hint{$x$ 与 $\dfrac{1}{2x}$ 均为负，先对 $-x$ 与 $-\dfrac{1}{2x}$ 用基本不等式，再把符号还原。与题 8 对照：只多一步转正。}
\end{problem}

\begin{problem}{题 10 （配凑定值）}
求下列函数的最小值：

(1) $y=x+\dfrac{1}{x-1}$（$x>1$）；\quad
(2) $y=\dfrac{x^2}{x-1}$（$x>1$）。
\hint{(1) 谁与 $\dfrac{1}{x-1}$ 相乘是定值？把 $x$ 写成 $(x-1)+1$；(2) 用 $x-1$ 改写分子：$x^2=(x-1)^2+2(x-1)+1$，逐项除开。}
\end{problem}

\begin{problem}{题 11 \starred（取倒数）}
求 $y=\dfrac{x-1}{x^2}$（$x>1$）的最大值。
\hint{它是题 10(2) 的倒数。$x>1$ 时 $y>0$，正数中较大者倒数较小，可直接引用 10(2) 的结论。}
\end{problem}

\begin{problem}{题 12 （和定积最大）}
求下列函数的最大值：

(1) $y=x(1-x)$（$0<x<1$）；\quad
(2) $y=x(1-2x)$（$0<x<\tfrac12$）。
\hint{(1) 两因子之和是多少？(2) 直接相加不是定值，先配系数：$y=\tfrac12\cdot 2x(1-2x)$，再看两因子之和。}
\end{problem}

\begin{problem}{题 13 （等号成立条件辨析）}
判断下列说法的真假，并说明理由：

A. $y=\sqrt{x^2+1}+\dfrac{4}{\sqrt{x^2+1}}$ 的最小值是 $4$；

B. 不等式 $\sqrt{x^2+1}+\dfrac{4}{\sqrt{x^2+1}}\ge 4$ 恒成立；

C. 对任意 $x>0$，$x^2+\dfrac{1}{x}\ge 2\sqrt{x}$ 恒成立；

D. 因为 $x=1$ 时 $x^2=\dfrac1x$（等号条件成立），所以 $y=x^2+\dfrac{1}{x}$（$x>0$）的最小值是 $y(1)=2$。
\hint{A、B 设 $t=\sqrt{x^2+1}$（注意 $t\ge 1$），等号要求 $t=2$，$x$ 取得到吗？C 与 D 形似而实异：C 只断言「$\ge$」，右边不必是定值；D 断言「最小值」，而 $2\sqrt{x^2\cdot\frac1x}=2\sqrt{x}$ 不是定值。取 $x=0.8$ 计算 $y$。}
\end{problem}

\clearpage

% ===========================================================================
%  板块五 · 含条件的最值（卡片 5 + 题 14–17）
% ===========================================================================
\section{含条件的最值}

\begin{strand}
\keyword{本板块主线}：用条件中的定值「1」沟通条件与目标。
\end{strand}

\block{知识精讲：三种处理方式}

\textbf{其一：乘「1」法（代换法）。}条件给出「和为 1」（或可化为和为定值），目标为倒数和时，把目标乘上这个「1」：
\[
a+b=1:\quad \frac1a+\frac1b=\left(\frac1a+\frac1b\right)(a+b)=2+\frac ba+\frac ab\ge 2+2=4.
\]
展开后常数项保留，互倒对用基本不等式处理。等号：$a=b=\tfrac12$，代回检验。（题 7 已用过三元版本。此法为何有效，见拓展 C：它在配平次数。）

\textbf{其二：和积互化（解不等式法）。}条件把 $a+b$ 与 $ab$ 联系在一起（如 $ab=a+b+3$）时，用 $a+b\ge 2\sqrt{ab}$ 或 $ab\le\left(\dfrac{a+b}2\right)^2$ 把其中一个换成另一个，得到单变量的不等式，解出范围。

\textbf{其三：双换元。}目标中两个分母各是 $x$、$y$ 的组合、无法各自化简时，设两个分母为新变量 $m$、$n$，反解出 $x$、$y$，条件与目标往往同时化简（题 17）。

\textbf{识别变形。}出题时常把「和为定值」隐藏起来：$\dfrac1a+\dfrac{2}{1-a}$ 即 $a+b=1$ 时的 $\dfrac1a+\dfrac2b$（令 $b=1-a$）；$\dfrac{1}{a+1}+\dfrac{1}{b+1}$ 中 $(a+1)+(b+1)=3$ 也是定值。先找出定值，再选方法。

\begin{problem}{题 14 （「1」的代换）}
已知 $a,b>0$ 且 $a+b=1$。

(1) 求 $\dfrac1a+\dfrac1b$ 的最小值；

(2) 求 $\dfrac1a+\dfrac2b$ 的最小值；

(3) 求 $\dfrac1a+\dfrac{2}{1-a}$（$0<a<1$）的最小值；

(4) \starred\ 求 $\dfrac{1}{a+1}+\dfrac{1}{b+1}$ 的最小值。
\hint{(2) 仍乘 $(a+b)$，展开后不互倒的两项同样可用基本不等式；(3) 先辨认它与哪一小问相同；(4) 令 $u=a+1$、$v=b+1$，$u+v$ 是定值吗？}
\end{problem}

\begin{problem}{题 15 \starred（和与积的互化）}
已知 $a,b>0$，且 $ab=a+b+3$，分别求 $ab$ 与 $a+b$ 的取值范围。
\hint{求 $ab$：用 $a+b\ge 2\sqrt{ab}$ 替换条件中的 $a+b$，设 $t=\sqrt{ab}$，移项后因式分解，注意其中一个因式恒正。求 $a+b$：改用「积 $\le$ 半和的平方」。端点能否取到，需回代检验（$a=b=3$）。}
\end{problem}

\begin{problem}{题 16 （乘「1」法）}
已知 $x,y>0$，且 $\dfrac1x+\dfrac9y=1$，求 $x+y$ 的最小值。
\hint{条件本身就是「1」，用 $x+y$ 乘它。与题 14(2) 对照：一处是「和为 1」，一处是「倒数和为 1」，方向相反。}
\end{problem}

\begin{problem}{题 17 \starred\starred（双换元）}
已知实数 $x,y$ 满足 $x>y>0$，且 $x+y\le 2$，求 $\dfrac{2}{x+3y}+\dfrac{1}{x-y}$ 的最小值。（书例 1.1-6，解答由讲义补全）
\hint{设 $m=x+3y$，$n=x-y$，先反解 $x$、$y$，看条件 $x+y\le 2$ 变成什么。目标 $\dfrac2m+\dfrac1n$ 没有「等于定值」的条件时，先乘 $(m+n)$ 再除回去，注意不等号的方向。}
\end{problem}

\clearpage

% ===========================================================================
%  课后自主练习（练习 1–13）
%  自总结版：在 main.tex 注释掉本文件的 \input 即可整节隐藏。
%  练习不设对应 solution 框，简解集中在 91-homework-solutions。
% ===========================================================================
\section{课后自主练习}


\begin{problem}{练习 1}
若 $0<a<b$ 且 $a+b=1$，则 $\dfrac12,\ a,\ 2ab,\ a^2+b^2$ 中最大的是哪个？
\hint{$a^2+b^2$ 与 $2ab$ 的大小是 \starref{}；再用 $a^2+b^2=(a+b)^2-2ab$ 与 $\dfrac12$ 比较。}
\end{problem}

\begin{problem}{练习 2}
已知 $a,b>0$。「$a+b\le 4$」能推出「$ab\le 4$」吗？反过来，「$ab\le 4$」能推出「$a+b\le 4$」吗？能推出的给出证明，不能的举反例。
\hint{正向用「积 $\le$ 半和的平方」；反向考虑乘积很小但相差悬殊的一对数。}
\end{problem}

\begin{problem}{练习 3 （多选）}
下列不等式中，\textbf{不一定}成立的是：

A. $x+\dfrac1x\ge 2$；\quad
B. $\sqrt{x^2+2}\ge\sqrt2$；\quad
C. $x^2+\dfrac{1}{x^2}\ge\sqrt2$（$x\ne0$）；\quad
D. $2-3x-\dfrac4x\ge 2$（$x\ne 0$）。
\hint{逐项检查 $x$ 是否允许取负；允许则代负数检验，只允许正数则检查等号。}
\end{problem}

\begin{problem}{练习 4}
设 $x,y>0$，求 $(x+y)\left(\dfrac1x+\dfrac4y\right)$ 的最小值。
\hint{卡片 3 旁注所指：先展开，常数保留，对剩余两项用基本不等式。}
\end{problem}

\begin{problem}{练习 5 \starred}
若 $-4<x<1$，求 $y=\dfrac{x^2-2x+2}{2x-2}$ 的最大值。
\hint{令 $t=x-1\in(-5,0)$，分子化为 $t^2+1$；负变量先转正（题 9），再确认取等点是否在区间内。}
\end{problem}

\begin{problem}{练习 6 \starred}
已知 $a,b>0$，若 $\dfrac2a+\dfrac1b\ge\dfrac{m}{2a+b}$ 恒成立，求 $m$ 的最大值。
\hint{分离 $m$：$m\le(2a+b)\left(\dfrac2a+\dfrac1b\right)$ 恒成立，等价于 $m\le$ 右边的最小值。展开右边。}
\end{problem}

\begin{problem}{练习 7}
已知 $x,y>0$ 且 $\dfrac x3+\dfrac y4=1$，求 $xy$ 的最大值。
\hint{和为定值、求积的最大值，只是这里的「和」是 $\dfrac x3+\dfrac y4$。}
\end{problem}

\begin{problem}{练习 8 \starred}
已知正数 $x,y$ 满足 $x+y=1$，求 $\dfrac{4}{x+2}+\dfrac{1}{y+1}$ 的最小值。
\hint{题 14(4) 的推广：$(x+2)+(y+1)=4$ 为定值，乘 $1=\dfrac{(x+2)+(y+1)}{4}$。}
\end{problem}

\begin{problem}{练习 9}
已知 $x>\dfrac12$，求 $f(x)=8x+\dfrac{1}{2x-1}$ 的最小值。（书例 1.2-1）
\hint{分母 $2x-1$ 不动，改写 $8x=4(2x-1)+4$。与题 10 同型。}
\end{problem}

\begin{problem}{练习 10}
若正数 $x,y$ 满足 $x+3y=5xy$，求 $3x+4y$ 的最小值。（书例 1.3-4）
\hint{条件两边同除以 $5xy$，构造出「1」，再按题 16 处理。}
\end{problem}

\begin{problem}{练习 11 \starred（多选）}
若实数 $x,y$ 满足 $x^2+y^2-xy=1$，则（\quad）。（书例 1.2-2）

A. $x+y<1$；\quad
B. $x+y\ge -2$；\quad
C. $x^2+y^2\le 2$；\quad
D. $x^2+y^2\ge 1$。
\hint{错误选项用特殊值排除（试 $x=y$ 与 $x=-y$）；正确选项一个用 $xy\le\dfrac{x^2+y^2}{2}$，一个把条件配方为 $\dfrac{(x+y)^2}{4}+\dfrac{3(x-y)^2}{4}$。$x,y$ 可负，可用的都是 \starref{} 一族（对一切实数成立）。}
\end{problem}

\begin{problem}{练习 12 \starred\starred}
已知 $5x^2y^2+y^4=1$（$x,y\in\mathbb{R}$），求 $x^2+y^2$ 的最小值。（书例 1.2-3）
\hint{左边分解为 $(5x^2+y^2)\cdot y^2$；两边乘 $4$，对 $5x^2+y^2$ 与 $4y^2$ 用「积 $\le$ 半和的平方」，半和恰为 $\dfrac52(x^2+y^2)$。}
\end{problem}

\begin{problem}{练习 13 \starred（选择）}
设正实数 $x,y,z$ 满足 $x^2-3xy+4y^2-z=0$，则当 $\dfrac{z}{xy}$ 取得最小值时，$x+2y-z$ 的最大值为（\quad）。

A. $0$；\quad
B. $\dfrac98$；\quad
C. $2$；\quad
D. $\dfrac94$。
\hint{先用条件消去 $z$；$\dfrac z{xy}$ 取最小值时 $x$、$y$ 的比例随之确定，回代后对 $y$ 配方。}
\end{problem}

\clearpage
     % 课后练习：做「自总结版」时整行注释
% ===========================================================================
%  拓展阅读（选读）· 拓展 A–D
%  可在 main.tex 注释掉本文件的 \input 整节隐藏。
% ===========================================================================
\section{拓展阅读（选读）}

\subsection*{拓展 A　平均数链}
\addcontentsline{toc}{subsection}{拓展 A　平均数链}

两个正数 $a,b$ 有四个常用的平均数：

\begin{center}
\begin{tabularx}{\linewidth}{@{}l l l L@{}}
\toprule
\multicolumn{1}{l}{\color{indigo}名称} &
\multicolumn{1}{l}{\color{indigo}记号} &
\multicolumn{1}{l}{\color{indigo}表达式} &
\multicolumn{1}{l}{\color{indigo}常见场景} \\
\midrule
调和平均 & $H$ & $\dfrac{2}{\frac1a+\frac1b}=\dfrac{2ab}{a+b}$ & 往返两段速度的平均速度 \\
几何平均 & $G$ & $\sqrt{ab}$ & 连续打折、增长率 \\
算术平均 & $A$ & $\dfrac{a+b}{2}$ & 通常所说的「平均」 \\
平方平均 & $Q$ & $\sqrt{\dfrac{a^2+b^2}{2}}$ & 物理中的方均根 \\
\bottomrule
\end{tabularx}
\end{center}

它们构成\keyword{平均数链}：
\[
H\le G\le A\le Q,\quad\text{每个等号都当且仅当 } a=b.
\]

\textbf{在一张图中同时呈现四个平均数。}（书 1.2.3）在板块二的半圆上添加几条线：直径 $BC$ 上取点 $D$，$BD=a$，$DC=b$，圆心 $O$；过 $D$ 的垂线交半圆于 $A$（板块二已证 $AD=\sqrt{ab}=G$）；取半圆弧中点 $E$（$OE\perp BC$，$OE=OA=\dfrac{a+b}2=A$）；再过 $D$ 作 $DF\perp AO$，垂足 $F$。

\begin{center}
\begin{tikzpicture}[scale=1.05, line cap=round]
  \coordinate (B) at (0,0);
  \coordinate (C) at (5,0);
  \coordinate (O) at (2.5,0);
  \coordinate (D) at (1.4,0);
  \coordinate (A) at (1.4,2.2449);
  \coordinate (E) at (2.5,2.5);
  \coordinate (F) at (2.2870,0.4347);
  \draw[inkiron] (B) -- (C);
  \draw[inkiron] (C) arc (0:180:2.5);
  \draw[selvage] (A) -- (O);
  \draw[weld, very thick] (A) -- (F) node[pos=0.78, left, weld]{\small $H$};
  \draw[madder, very thick] (D) -- (A) node[midway, left, madder]{\small $G$};
  \draw[indigo, very thick] (O) -- (E) node[midway, right, indigo]{\small $A$};
  \draw[selvage, very thick] (D) -- (E) node[pos=0.22, right, selvage]{\small $Q$};
  \draw[inkiron, thin] (D) -- (F);
  \foreach \p/\n/\where in {B/B/below, C/C/below, O/O/below, D/D/below, A/A/above, E/E/above, F/F/right}{
    \fill[inkiron] (\p) circle (1.3pt);
    \node[\where, inkiron] at (\p) {\small $\n$};
  }
  \draw[weld, |-|] (0,-0.55) -- node[below, weld]{\small $a$} (1.4,-0.55);
  \draw[weld, |-|] (1.4,-0.55) -- node[below, weld]{\small $b$} (5,-0.55);
\end{tikzpicture}
\end{center}

由于 $AD\perp BC$ 且 $DO$ 在 $BC$ 上，$\triangle ADO$ 在 $D$ 处成直角，对它用射影定理（0.3）：
\[
AD^2=AF\cdot AO\ \Longrightarrow\ AF=\frac{AD^2}{AO}=\frac{ab}{\frac{a+b}{2}}=\frac{2ab}{a+b}=H;
\]
再在 $\mathrm{Rt}\triangle DOE$（直角在 $O$）中用勾股定理：
\[
DE=\sqrt{OD^2+OE^2}=\sqrt{\left(\frac{b-a}{2}\right)^2+\left(\frac{a+b}{2}\right)^2}=\sqrt{\frac{a^2+b^2}{2}}=Q.
\]
读图，三次使用「直角三角形中斜边不小于直角边」：
% \underbrace 在容器数学字体下会丢字形（见 latex-pitfalls），改用 array 双行标注
\[
\begin{array}{ccccccc}
DE &\ \ge\ & OE=OA &\ \ge\ & AD &\ \ge\ & AF, \\[1pt]
Q  &       & A     &       & G  &       & H
\end{array}
\]
（依次在 $\mathrm{Rt}\triangle DOE$、$\mathrm{Rt}\triangle ADO$、$\mathrm{Rt}\triangle AFD$ 中。）一张图给出整条链。$D$ 移动到圆心时四条线段重合，对应 $a=b$ 时四个平均数相等。

\textbf{代数复核}：中间 $G\le A$ 是基本不等式；左边 $H\le G$ 是把基本不等式用于 $\dfrac1a,\dfrac1b$ 后取倒数；右边 $A\le Q$ 两边平方后即 \starref{}。三段证明请自行补全。

\textbf{应用：不等臂天平。}天平两臂长 $l_1\ne l_2$。先把 10g 砝码放左盘、药放右盘称出一份，再把砝码放右盘、药放左盘称出一份。由杠杆平衡（力 $\times$ 臂长相等）：第一份实重 $m_1=\dfrac{10\,l_1}{l_2}$，第二份 $m_2=\dfrac{10\,l_2}{l_1}$。于是 $\sqrt{m_1m_2}=10$ 为定值，而
\[
m_1+m_2=10\left(\frac{l_1}{l_2}+\frac{l_2}{l_1}\right)\ge 20,
\]
$l_1\ne l_2$ 时严格大于 20：两次称量之和总是超过 20g。两次称量的算术平均偏大，与真实质量对应的是几何平均。

\textbf{推广到 $n$ 个正数}（记住形式即可）：
\[
\sqrt{\frac{x_1^2+\cdots+x_n^2}{n}}\ \ge\ \frac{x_1+\cdots+x_n}{n}\ \ge\ \sqrt[n]{x_1\cdots x_n}\ \ge\ \frac{n}{\frac1{x_1}+\cdots+\frac1{x_n}}.
\]
其中三元的「和 $\ge$ 3 倍立方根积」最常用：$p+q+r\ge 3\sqrt[3]{pqr}$（$p,q,r>0$）。示例：

\textbf{例}　$a,b,c>0$，$abc=1$，证明 $(a+b)^3+(b+c)^3+(c+a)^3\ge 24$。

先对三个立方用三元均值，再对每个括号用基本不等式：
\[
(a+b)^3+(b+c)^3+(c+a)^3\ge 3(a+b)(b+c)(c+a)\ge 3\cdot 2\sqrt{ab}\cdot 2\sqrt{bc}\cdot 2\sqrt{ca}=24\sqrt{(abc)^2}=24,
\]
当且仅当 $a=b=c=1$ 时取等。两层放缩、取等条件一致，正是板块三的方法。

\newpage
\subsection*{拓展 B　用一个函数统一四个平均数 \starred\starred}
\addcontentsline{toc}{subsection}{拓展 B　用一个函数统一四个平均数}

设 $0<a\le b$，考察
\[
f(x)=\frac{b-x}{x-a}=-1+\frac{b-a}{x-a}\quad(x>a).
\]
先确定单调性：$x$ 增大时分母 $x-a$ 增大，$\dfrac{b-a}{x-a}$ 减小，故 $f$ 在 $(a,+\infty)$ 上递减。逐一代入（均可直接验证）：
\[
f(A)=1,\qquad f(H)=\frac ba,\qquad f(G)=\sqrt{\frac ba},\qquad f(Q)=\frac{Q+a}{b+Q}.
\]
由 $\dfrac ba\ge\sqrt{\dfrac ba}\ge 1\ge\dfrac{Q+a}{b+Q}$（末项成立是因为 $Q\le b$，分子不超过分母；请自行补证 $a\le Q\le b$），且 $f$ 递减，函数值大者自变量小，即得
\[
H\le G\le A\le Q.
\]
四个平均数由一个减函数的四个函数值统一刻画。

\newpage
\subsection*{拓展 C　条件最值问题的生成：齐次化的视角}
\addcontentsline{toc}{subsection}{拓展 C　条件最值问题的生成：齐次化的视角}

平均数链上，知道一个平均数的值，其余三个的范围随之确定（「知一求三」）。例如 $\sqrt{ab}=1$（即 $ab=1$）立即给出 $\dfrac{a+b}{2}\ge 1$。许多条件最值题，都是把这个基础模型逐步变形、隐藏得到的：

\textbf{第 0 步（基础型）}　$a,b>0$，$ab=1$，则 $a+b\ge 2$。

\textbf{第 1 步（改写变量）}　把基础型「$uv=9$，求 $u+v$ 最小」中的变量写作 $u=a+1$、$v=2b+1$，题面变为「$(a+1)(2b+1)=9$（$a,b>0$），求 $a+2b$ 的最小值」。识别出 $u+v=a+2b+2$ 后本质未变：$(a+1)+(2b+1)\ge 2\sqrt9=6$，故 $a+2b\ge 4$（$a=2,\ b=1$ 取等）。（乘积由 1 调为 9，是为使取等点落在正数范围内。）

\textbf{第 2 步（展开隐藏结构）}　把条件展开：$2ab+a+2b=8$，求 $a+2b$ 的最小值。与第 1 步是同一道题，但需要自行把 $(a+1)(2b+1)=9$ 因式分解还原。此步是解这类题的关键功夫。

\textbf{第 3 步（调整系数）}　条件改为 $2ab+a+2b=2$（即 $(a+1)(2b+1)=3$），求 $2a+2b$ 的最小值。目标系数与条件不齐时配凑：两边乘 2 得 $(2a+2)(2b+1)=6$，而 $(2a+2)+(2b+1)\ge 2\sqrt6$，故 $2a+2b\ge 2\sqrt6-3$。

\textbf{第 4 步（升为二次）}　条件改成二次式：$a^2+3ab+2b^2=1$，求 $5a+7b$ 的最小值。先因式分解 $a^2+3ab+2b^2=(a+b)(a+2b)$，积为定值；再拆目标 $5a+7b=3(a+b)+2(a+2b)\ge 2\sqrt{3\cdot2\cdot1}=2\sqrt6$（$a=b=\dfrac{1}{\sqrt6}$ 取等）。

\textbf{这些变形共同的引擎：齐次。}每一项次数都相同的式子称为\keyword{齐次式}：$a+b$ 是 1 次，$ab$、$a^2+3ab+2b^2$ 是 2 次，$\dfrac ba+\dfrac ab$ 是 0 次。基本不等式 $a+b\ge 2\sqrt{ab}$ 两边都是 1 次——\textbf{它只比较同次的量}。由此可以解释两件事：
\begin{itemize}
  \item \textbf{0 次齐次式只依赖比值。}$2+\dfrac ba+\dfrac ab$ 中设 $t=\dfrac ba$，即化为 $2+t+\dfrac1t$，双变量化为单变量。这是「乘 1 展开后能继续处理」的根本原因。
  \item \textbf{乘「1」法等于配平次数。}$\dfrac1a+\dfrac1b$ 是 $-1$ 次，乘上 1 次的 $(a+b)$ 恰为 0 次。条件（$a+b=1$、$ab=1$、$\dfrac1x+\dfrac9y=1$ 等）的作用，是提供一个可自由乘除的「1」，用来补齐目标式的次数。换元、配凑、乘 1，做的是同一件事：\keyword{把式子修成齐次，使基本不等式可用}。
\end{itemize}

书中的一道示范，分子各项次数不一也能配齐：

\textbf{例}　正实数 $a,b$ 满足 $a+2b=2$，求 $\dfrac{1+4a+3b}{ab}$ 的最小值。

分母 $ab$ 是 2 次；分子中 $1$ 是 0 次、$4a+3b$ 是 1 次，用条件逐项补齐到 2 次：$1=\left(\dfrac{a+2b}{2}\right)^2$，$4a+3b=(4a+3b)\cdot\dfrac{a+2b}{2}$。于是
\begin{align*}
\frac{1+4a+3b}{ab}
  &=\frac{\left(\frac{a+2b}{2}\right)^2+(4a+3b)\cdot\frac{a+2b}{2}}{ab}
   =\frac{9a^2+26ab+16b^2}{4ab}\\
  &=\frac{26}{4}+\frac{9a^2+16b^2}{4ab}
   \ge\frac{26}{4}+\frac{2\sqrt{144}\,ab}{4ab}=\frac{25}{2},
\end{align*}
当且仅当 $9a^2=16b^2$ 即 $3a=4b$，配合 $a+2b=2$ 得 $a=\dfrac45$，$b=\dfrac35$ 时取等。

了解这条生成路径后，面对陌生的条件最值题可反向操作，先问两个问题：条件能否因式分解还原为基础型？目标与条件各是几次、相差几次？

\newpage
\subsection*{拓展 D　对称性：可以用来猜，不能用来证}
\addcontentsline{toc}{subsection}{拓展 D　对称性：可以用来猜，不能用来证}

\textbf{轮换对称式。}把式中字母按一定顺序轮换（如 $x\to y$，$y\to x$；三元时 $x\to y\to z\to x$）后式子不变，称其为轮换对称式。（书 1.3.2）基本不等式一族的取等条件都是「各元相等」，因此面对对称的问题，很自然会猜测取等点在 $x=y$。这一猜测何时可靠？先看同一道题的三种解法：

\textbf{同一题三解}　实数 $x,y$ 满足 $x^2+y^2+xy=1$，求 $x+y$ 的最大值。（书例 1.3-1）

\textbf{解法一（配凑）}：$x^2+y^2+xy=(x+y)^2-xy$，而 $xy\le\left(\dfrac{x+y}{2}\right)^2$ 对一切实数成立（卡片 1 易错提醒），故
\[
1=(x+y)^2-xy\ \ge\ (x+y)^2-\frac{(x+y)^2}{4}=\frac34(x+y)^2,
\]
得 $(x+y)^2\le\dfrac43$，$x+y\le\dfrac{2\sqrt3}{3}$，等号在 $x=y=\dfrac{\sqrt3}{3}$（代回约束验证成立）。

\textbf{解法二（对称猜测）}：题目关于 $x,y$ 对称，令 $x=y$：$3x^2=1$，$x+y=2x=\dfrac{2\sqrt3}{3}$。此法只算出了对称点处的值，并未排除其他点更大，\textbf{只能作为猜测}，需用解法一或解法三坐实。

\textbf{解法三（判别式法）}：设 $k=x+y$，以 $x=k-y$ 代入约束得 $y^2-ky+k^2-1=0$。$y$ 为实数，方程必须有实数解，故判别式 $\Delta=k^2-4(k^2-1)\ge 0$（一元二次方程有实数解当且仅当 $\Delta=b^2-4ac\ge 0$，初三内容，此处直接使用），解得 $k^2\le\dfrac43$，结论相同。

同型题：$4x^2+y^2+xy=1$ 求 $2x+y$ 的最大值——把 $2x$ 视为一个整体，三种解法均适用，答案 $\dfrac{2\sqrt{10}}{5}$。（书例 1.2-5/1.3-3）

\textbf{对称猜测失效的例子。}$a,b>0$，$ab=1$，求 $\dfrac{1}{2a}+\dfrac{1}{2b}+\dfrac{8}{a+b}$ 的最小值。

问题关于 $a,b$ 完全对称，猜 $a=b=1$，得 $\dfrac12+\dfrac12+4=5$。\textbf{这是错误答案。}正确解法：用 $ab=1$ 合并前两项，
\[
\frac{1}{2a}+\frac{1}{2b}=\frac{a+b}{2ab}=\frac{a+b}{2},
\]
设 $t=a+b\ge 2\sqrt{ab}=2$，目标化为 $\dfrac t2+\dfrac 8t\ge 2\sqrt{\dfrac t2\cdot\dfrac8t}=4$，等号要求 $t=4$，在 $t\ge 2$ 范围内可以取到（$a+b=4$ 且 $ab=1$，解得 $a=2\pm\sqrt3$）。最小值为 $4$，且取等点并不对称。

失效原因：换元后 $g(t)=\dfrac t2+\dfrac8t$ 的最小值点在 $t=4$，而对称点 $a=b$ 对应的是端点 $t=2$。对称只说明「对称点是一个特殊位置」，不保证最值在该处。因此：

\begin{strand}
\keyword{齐次决定结构，对称提示取等点，证明必须由不等式本身完成。}用对称猜出取等点后，务必用配凑或判别式法验证；使用判别式法求最值时，还要验证端点能否取到、变量是否有附加限制（如 $\sqrt{y}\ge 0$ 一类隐含条件，参见书例 1.3-2）。
\end{strand}

\clearpage
   % 拓展阅读（选读）：可整行注释

\clearpage
% ---- 第二部分：集中解析（与题目区物理分开）----
% ===========================================================================
%  第二部分 · 参考答案与详细解析（题 1–17，四段式，与题目区一一对应）
% ===========================================================================
\section*{第二部分\quad 参考答案与详细解析}

\begin{strand}
每题按「怎么想 → 解 → 检验 → 方法小结」展开。先看「怎么想」核对思路方向，再看解答核对细节。
\end{strand}

\begin{solution}{题 1 · 解析（作差与公式的选择）}
\think (1) 是「互倒之和」的标准形式。(2) 的关键在 $a$ 可取负值：基本不等式的正数前提不满足，应退回作差配方。(3) 换成 2 后，配方所得的常数项还剩多少？

\solve (1) 因为 $a,b>0$，所以 $\dfrac ab>0,\ \dfrac ba>0$，且 $\dfrac ab\cdot\dfrac ba=1$ 为定值。由基本不等式
\[
\frac ab+\frac ba\ge 2\sqrt{\frac ab\cdot\frac ba}=2,
\]
当且仅当 $\dfrac ab=\dfrac ba$，即 $a=b$ 时取等号。故 $\dfrac ab+\dfrac ba\ge 2$（$a=b$ 时相等）。

(2) 作差：$a^2+3-3a=\left(a-\dfrac32\right)^2+\dfrac34\ge\dfrac34>0$，对一切实数 $a$ 成立。故 $a^2+3>3a$，且等号永远取不到。

(3) $a^2+2-2a=(a-1)^2\ge 0$，即 $a^2+2\ge 2a$，$a=1$ 时相等。常数换成 2 后恰好配成完全平方，「严格大于」变为「可以相等」。

\checkup (2) 取 $a=\dfrac32$（差的最小值点）：左边 $=\dfrac{21}4$，右边 $=\dfrac{18}4$，差为 $\dfrac34$，与配方结果一致。(3) 取 $a=1$：两边均为 3。

\insight{变量可能取负时先作差配方；确认各项为正后才可使用基本不等式。配方后的常数项决定「严格大于」还是「可以相等」。}
\end{solution}

\begin{solution}{题 2 · 解析（四数排序）}
\think 先猜排序：$a$ 最小、$b$ 最大、两个平均数居中、算术平均在几何平均之上（引例已给出方向）。三个空隙逐段作差或提公因式。

\solve 设 $0<a<b$。

① $\sqrt{ab}-a=\sqrt a(\sqrt b-\sqrt a)>0$（因 $\sqrt b>\sqrt a>0$），故 $a<\sqrt{ab}$。

② 由基本不等式且 $a\ne b$：$\sqrt{ab}<\dfrac{a+b}{2}$（等号条件不满足，取严格不等号）。

③ $b-\dfrac{a+b}{2}=\dfrac{b-a}{2}>0$，故 $\dfrac{a+b}{2}<b$。

综上：$a<\sqrt{ab}<\dfrac{a+b}{2}<b$。

\checkup 取 $a=1,b=9$：$1<3<5<9$，成立。

\insight{两个平均数总被夹在原来两数之间。对照板块二的图：$\sqrt{ab}$（半弦）与 $\frac{a+b}2$（半径）都比整条直径短、比较短的一段长。}
\end{solution}

\begin{solution}{题 3 \starred\ · 解析（两步放缩）}
\think 一次基本不等式不够：若直接写 $a^2+\dfrac2a\ge 2\sqrt{2a}$，右边含 $a$，不是常数（与题 13 D 同类错误）。改走两步，每一步都朝「消去平方、凑出互倒」进行。

\solve 因为 $a>0$：

第一步：$a^2+1\ge 2a$（即 \starref{}），所以
\[
a^2+\frac2a=(a^2+1)+\frac2a-1\ge 2a+\frac2a-1.
\]
第二步：$2a+\dfrac2a=2\left(a+\dfrac1a\right)\ge 2\cdot 2=4$。

由传递性（0.1 旁注）：$a^2+\dfrac2a\ge 4-1=3$。

取等条件核对：第一步要求 $a=1$，第二步要求 $a=\dfrac1a$ 即 $a=1$，两步一致。故 $a^2+\dfrac2a\ge 3$，当且仅当 $a=1$ 时取等号。

\checkup $a=1$：$1+2=3$；$a=2$：$4+1=5>3$。

\insight{链式放缩中，各步取等条件必须在同一点成立。若取等点不同，合成的「$\ge$」仍然正确，但端点值取不到，得到的只是下界而非最值。}
\end{solution}

\begin{solution}{题 4 · 解析（下界与上界）}
\think $m$ 是配凑互倒的典型：$a>2$ 时 $a-2>0$，拆 $a=(a-2)+2$。$n$ 是有上界的二次式：$b\ne 0$ 使它到不了 4。$m$ 的下界是 4 且可以取到，$n$ 的上界是 4 且取不到，两者由此分开。

\solve 因为 $a>2$，$a-2>0$：
\[
m=(a-2)+\frac{1}{a-2}+2\ge 2\sqrt{(a-2)\cdot\frac{1}{a-2}}+2=4,
\]
当且仅当 $a-2=\dfrac{1}{a-2}$，即 $a=3$ 时取等号，故 $m\ge 4$。

又 $b\ne 0$ 时 $b^2>0$，故 $n=4-b^2<4$。

所以 $n<4\le m$，即 $m>n$。

\checkup $a=3$：$m=4$；$b=0.1$：$n=3.99<4$。

\insight{分母确定后，把分母以外的部分改写为「分母 + 常数」。比较两个变量的大小，常分别求出一方的下界与另一方的上界，验证两条界线是否分开。}
\end{solution}

\begin{solution}{题 5 · 解析（半圆中的构造）}
\think 两条预备定理各司其职：圆周角定理提供直角三角形，射影定理给出 $\sqrt{ab}$。最后一步只需「斜边不小于直角边」。

\solve (1) $AB$ 为直径，$D$ 在圆上，由 0.3 圆周角定理，$\angle ADB=90^\circ$。于是 $\triangle ADB$ 是直角三角形，$DE$ 是斜边 $AB$ 上的高，垂足 $E$ 把斜边分成 $AE=a$、$EB=b$。由射影定理：
\[
DE^2=AE\cdot EB=ab,\qquad DE=\sqrt{ab}.
\]

(2) $DO$ 为半径，$DO=\dfrac{a+b}{2}$。在 $\mathrm{Rt}\triangle DEO$ 中（直角在 $E$），$DO$ 是斜边，$DE$ 是直角边，故
\[
DE\le DO,\qquad\text{即}\qquad \sqrt{ab}\le\frac{a+b}{2}.
\]

(3) 等号成立当且仅当三角形退化，即 $E$ 与 $O$ 重合。此时 $AE=EB$，$a=b$，$D$ 位于半圆最高点，$DE$ 即为一条半径。

\checkup $a=1,b=9$：半径 5，$OE=4$，$DE=3$，$3^2+4^2=5^2$，且 $3<5$，与勾股定理及结论一致。

\insight{半弦不过半径：一句话概括不等式本身与取等条件——半弦等于半径，当且仅当垂足与圆心重合。}

\begin{remark*}
延伸（不使用 0.3 的另一条路）：在过 $E$ 的垂线上取点 $D'$ 使 $ED'=\sqrt{ab}$，由勾股定理算得 $OD'^2=\left(\dfrac{b-a}{2}\right)^2+ab=\left(\dfrac{a+b}{2}\right)^2$，即 $OD'$ 等于半径，$D'$ 落在圆上，与 $D$ 为同一点。两条路线可互相印证。
\end{remark*}
\end{solution}

\begin{solution}{题 6 · 解析（逐项放缩相乘）}
\think 右边 $8x^3y^3=2\cdot 2\cdot 2\cdot x^3y^3$，提示三次基本不等式各出一个因子 2；三个括号分别对应 $x\cdot y$、$x^2\cdot y^2$、$x^3\cdot y^3$。

\solve 因为 $x,y>0$，下列三式两边均为正：
\[
x+y\ge 2\sqrt{xy},\qquad x^2+y^2\ge 2\sqrt{x^2y^2}=2xy,\qquad x^3+y^3\ge 2\sqrt{x^3y^3}.
\]
同向相乘：
\[
(x+y)(x^2+y^2)(x^3+y^3)\ge 8\sqrt{xy}\cdot xy\cdot\sqrt{x^3y^3}=8\,xy\cdot\sqrt{x^4y^4}=8x^3y^3.
\]
三处等号都当且仅当 $x=y$，可同时成立。故原不等式成立，$x=y$ 时取等。

\checkup $x=y=1$：左边 $=8$，右边 $=8$；$x=1,y=2$：左边 $=3\cdot5\cdot9=135$，右边 $=64$，不等式成立。

\insight{正数不等式可同向相乘，取等条件必须全体一致。右边为 $2^k\times$ 单项式时，可先考虑 $k$ 次基本不等式相乘。}
\end{solution}

\begin{solution}{题 7 \starred\ · 解析（用条件替换 1）}
\think 目标中的分子都是 1，条件给出 $1=a+b+c$。把 1 原地替换，每个分式裂为三项，整理出 3 个常数与 3 对互倒。

\solve 由 $a+b+c=1$：
\[
\frac1a+\frac1b+\frac1c=\frac{a+b+c}{a}+\frac{a+b+c}{b}+\frac{a+b+c}{c}
=3+\left(\frac ba+\frac ab\right)+\left(\frac ca+\frac ac\right)+\left(\frac cb+\frac bc\right).
\]
三对互倒之和各不小于 2（题 1(1)），故右边 $\ge 3+2+2+2=9$。三处等号分别要求 $a=b$、$a=c$、$b=c$，即 $a=b=c=\dfrac13$，可同时成立。

\checkup $a=b=c=\dfrac13$：左边 $=3+3+3=9$。

\insight{条件「和为 1」的标准用法：把目标中的 1 替换为这个和。「乘 1」与「换 1」实为同一操作，板块五反复使用。}
\end{solution}

\begin{solution}{题 8 · 解析（直接使用）}
\think 三查：一正（$x>0$，两项皆正）；二定（$4x\cdot\dfrac{12}x=48$ 为定值）；三相等（等号点是否在定义域内，算出来核对）。

\solve $x>0$ 时，
\[
y=4x+\frac{12}{x}\ge 2\sqrt{4x\cdot\frac{12}{x}}=2\sqrt{48}=8\sqrt3,
\]
当且仅当 $4x=\dfrac{12}{x}$，即 $x^2=3$，$x=\sqrt3$（取正）时取等号。最小值为 $8\sqrt3$（此时 $x=\sqrt3$）。

\checkup $x=\sqrt3$：$y=4\sqrt3+4\sqrt3=8\sqrt3$；$x=1$：$y=16>8\sqrt3\approx 13.86$。

\insight{积为定值 $P$ 时，和的最小值为 $2\sqrt{P}$，结果常含根号。报最小值必须同时给出取到最小值的位置。}
\end{solution}

\begin{solution}{题 9 · 解析（负变量先转正）}
\think $x<0$ 时两项均为负，不满足「一正」。对 $-x>0$、$-\dfrac{1}{2x}>0$ 使用基本不等式，再把符号还原；不等号方向随之反转，最小值问题变为最大值问题。

\solve $x<0$ 时，$-x>0$，$-\dfrac{1}{2x}>0$，且积 $(-x)\left(-\dfrac{1}{2x}\right)=\dfrac12$ 为定值：
\[
(-x)+\left(-\frac{1}{2x}\right)\ge 2\sqrt{\frac12}=\sqrt2,
\]
即 $x+\dfrac{1}{2x}\le-\sqrt2$，从而 $y=x+\dfrac{1}{2x}-2\le-\sqrt2-2$。

等号当且仅当 $-x=-\dfrac{1}{2x}$，即 $x^2=\dfrac12$，$x=-\dfrac{\sqrt2}{2}$（取负）。最大值为 $-2-\sqrt2$。

\checkup $x=-\dfrac{\sqrt2}2$：$x+\dfrac1{2x}=-\sqrt2$，$y=-\sqrt2-2$，与结论一致。

\insight{负区间求最值：先取相反数转正，用完基本不等式再把不等号方向还原。与题 8 相比只多一步转正。}
\end{solution}

\begin{solution}{题 10 · 解析（配凑定值）}
\think (1) $x$ 与 $\dfrac1{x-1}$ 之积不是定值，但 $(x-1)$ 与 $\dfrac{1}{x-1}$ 互为倒数，故把 $x$ 拆为 $(x-1)+1$。(2) 分母是 $x-1$，就用 $x-1$ 改写分子。

\solve (1) $x>1$ 时 $x-1>0$：
\[
y=(x-1)+\frac{1}{x-1}+1\ge 2\sqrt{(x-1)\cdot\frac{1}{x-1}}+1=3,
\]
当且仅当 $x-1=\dfrac1{x-1}$，即 $x=2$ 时取等。最小值 $3$。

(2) 分子 $x^2=(x-1)^2+2(x-1)+1$，故
\[
y=\frac{x^2}{x-1}=(x-1)+\frac{1}{x-1}+2\ge 2+2=4,
\]
当且仅当 $x=2$ 时取等。最小值 $4$。

\checkup (1) $x=2$：$y=3$。(2) $x=2$：$y=4$；$x=3$：$y=4.5>4$。

\insight{分母确定后，把式子改写为「分母 + 分母的倒数 + 常数」。配凑不改变函数本身，只改变观察角度。（(1) 是题 8 经变量平移得到的变形，参见拓展 C。）}
\end{solution}

\begin{solution}{题 11 \starred\ · 解析（取倒数）}
\think 分子一次、分母二次，直接处理不便；其倒数 $\dfrac{x^2}{x-1}$ 恰是题 10(2)。正数中较大者倒数较小，最小值 4 对应倒数的最大值 $\dfrac14$。

\solve $x>1$ 时 $x-1>0$，故 $y=\dfrac{x-1}{x^2}>0$。由题 10(2)：$\dfrac{x^2}{x-1}\ge 4$，两边取倒数（均为正，方向反转）：
\[
y=\frac{x-1}{x^2}\le\frac14,
\]
当且仅当 $x=2$ 时取等。最大值 $\dfrac14$。

\checkup $x=2$：$y=\dfrac{1}{4}$；$x=3$：$y=\dfrac29<\dfrac14$。

\insight{结构不便处理时可考察倒数。正数间取倒数反转大小关系：最小值的倒数即倒数的最大值，转换前必须确认符号为正。}
\end{solution}

\begin{solution}{题 12 · 解析（和定积最大）}
\think (1) $x+(1-x)=1$，和为定值。(2) $x+(1-2x)$ 不是定值，配系数使两因子之和消去 $x$：$2x+(1-2x)=1$。

\solve (1) $0<x<1$ 时 $x>0,\ 1-x>0$，两者之和为 1：
\[
y=x(1-x)\le\left(\frac{x+(1-x)}{2}\right)^2=\frac14,
\]
当且仅当 $x=1-x$，即 $x=\dfrac12$ 时取等。最大值 $\dfrac14$。

(2) $0<x<\dfrac12$ 时 $2x>0,\ 1-2x>0$：
\[
y=x(1-2x)=\frac12\cdot 2x(1-2x)\le\frac12\left(\frac{2x+(1-2x)}{2}\right)^2=\frac12\cdot\frac14=\frac18,
\]
当且仅当 $2x=1-2x$，即 $x=\dfrac14$ 时取等。最大值 $\dfrac18$。

\checkup (1) $x=\dfrac12$：$y=\dfrac14$。(2) $x=\dfrac14$：$y=\dfrac14\cdot\dfrac12=\dfrac18$。

\insight{和为定值 $S$ 时积的最大值为 $\frac{S^2}{4}$；和不是定值时先配系数凑成定值，配完在外面除回去。几何解释：周长一定的矩形中，正方形面积最大。}
\end{solution}

\begin{solution}{题 13 · 解析（等号成立条件辨析）}
\think 四个说法都围绕「三相等」。A、B 换元后检查等号点是否可达；C 只断言不等式恒成立，不要求右边是定值；D 把「等号条件成立处的函数值」当成了「最小值」，而右边 $2\sqrt x$ 不是定值，推理不成立。

\solve A（真）。设 $t=\sqrt{x^2+1}\ge 1$，$y=t+\dfrac4t\ge 2\sqrt4=4$，等号要求 $t=2$，即 $x^2+1=4$，$x=\pm\sqrt3$，可以取到。最小值确为 4。

B（真）。由 A 的推理，$y\ge 4$ 对一切实数 $x$ 成立。

C（真）。$x>0$ 时，$x^2>0,\ \dfrac1x>0$，由基本不等式 $x^2+\dfrac1x\ge 2\sqrt{x^2\cdot\dfrac1x}=2\sqrt x$，恒成立（等号在 $x=1$）。

D（假）。$x=1$ 时 $y=2$ 只是一个函数值。$2\sqrt x$ 不是定值，「$y\ge 2\sqrt x$」不能给出最小值。反例：$x=0.8$ 时 $y=0.64+1.25=1.89<2$，故 2 不是最小值。

\checkup A：$x=\sqrt3$ 时 $t=2$，$y=4$。D 的反例已代入验证。

\insight{「$\ge$」的右边是定值、且等号可以取到，二者齐备才能宣布最值。C 与 D 的对照说明：同一个不等式，「恒成立」为真而「最小值」为假，二者不可混同。}
\end{solution}

\begin{solution}{题 14 · 解析（「1」的代换）}
\think 四小问是同一方法的四种形态：(1)(2) 直接乘 1 展开；(3) 是 (2) 的变量替换形式；(4) 先设新变量，再乘新的定值和。

\solve $a,b>0,\ a+b=1$。

(1) $\dfrac1a+\dfrac1b=\left(\dfrac1a+\dfrac1b\right)(a+b)=2+\dfrac ba+\dfrac ab\ge 2+2=4$，当且仅当 $a=b=\dfrac12$ 取等。最小值 $4$。

(2) $\dfrac1a+\dfrac2b=\left(\dfrac1a+\dfrac2b\right)(a+b)=1+2+\dfrac ba+\dfrac{2a}{b}\ge 3+2\sqrt{\dfrac ba\cdot\dfrac{2a}b}=3+2\sqrt2$，当且仅当 $\dfrac ba=\dfrac{2a}b$，即 $b=\sqrt2 a$，配合 $a+b=1$ 得 $a=\sqrt2-1,\ b=2-\sqrt2$（均为正）。最小值 $3+2\sqrt2$。

(3) 令 $b=1-a\in(0,1)$，则 $\dfrac1a+\dfrac{2}{1-a}=\dfrac1a+\dfrac2b$ 且 $a+b=1$，与 (2) 相同。最小值 $3+2\sqrt2$。

(4) 令 $u=a+1,\ v=b+1$，则 $u,v>0$ 且 $u+v=3$：
\[
\frac1u+\frac1v=\frac13\left(\frac1u+\frac1v\right)(u+v)=\frac13\left(2+\frac vu+\frac uv\right)\ge\frac13(2+2)=\frac43,
\]
当且仅当 $u=v=\dfrac32$，即 $a=b=\dfrac12$ 取等。最小值 $\dfrac43$。

\checkup (2) $a=\sqrt2-1$：$\dfrac1a=\sqrt2+1$，$\dfrac2b=\dfrac{2}{2-\sqrt2}=2+\sqrt2$，两项之和 $3+2\sqrt2$。(4) $a=b=\dfrac12$：$\dfrac{1}{3/2}\times2=\dfrac43$。

\insight{条件和为定值 $S$、目标为倒数和时，乘上 $\frac{\text{和}}{S}$ 这个「1」，常数项保留、互倒对用基本不等式。常见的两种伪装是变量平移（(3)）与整体换元（(4)），先找出定值和即可还原。}
\end{solution}

\begin{solution}{题 15 \starred\ · 解析（和与积的互化）}
\think 条件 $ab=a+b+3$ 中和与积相互纠缠，都不是定值。方法是把其中一个换成另一个：求 $ab$ 时用 $a+b\ge 2\sqrt{ab}$ 消去 $a+b$；求 $a+b$ 时用 $ab\le\left(\frac{a+b}2\right)^2$ 消去 $ab$，各得一个单变量的二次不等式。

\solve 求 $ab$：由 $a+b\ge 2\sqrt{ab}$，得 $ab=a+b+3\ge 2\sqrt{ab}+3$。设 $t=\sqrt{ab}>0$：
\[
t^2-2t-3\ge 0\ \Longrightarrow\ (t-3)(t+1)\ge 0.
\]
因为 $t>0$，因式 $t+1$ 恒正，故必有 $t-3\ge 0$，即 $t\ge 3$，$ab=t^2\ge 9$。

求 $a+b$：由 $ab\le\left(\dfrac{a+b}2\right)^2$，设 $s=a+b>0$：
\[
s+3\le\frac{s^2}{4}\ \Longrightarrow\ s^2-4s-12\ge 0\ \Longrightarrow\ (s-6)(s+2)\ge 0,
\]
同理 $s+2>0$ 恒正，故 $s\ge 6$。

端点检验：$a=b=3$ 时 $ab=9$，$a+b+3=9$，条件满足，两个端点同时取到。故 $ab\ge 9$，$a+b\ge 6$（$a=b=3$ 时取等）。

\checkup 另取满足条件的 $a=5$：$5b=5+b+3$ 得 $b=2$，此时 $ab=10\ge 9$，$a+b=7\ge 6$。

\insight{和与积纠缠时，用基本不等式把双变量压成单变量，因式分解解二次不等式（恒正的因式直接约去）。端点必须回代检验能否取到，这是「三相等」在此类问题中的形式。}
\end{solution}

\begin{solution}{题 16 · 解析（乘「1」法）}
\think 条件 $\dfrac1x+\dfrac9y=1$ 本身就是「1」。用 $x+y$ 乘它，展开后常数与互倒对分别处理。与题 14 对照：一处「和为 1」乘给倒数和，一处「倒数和为 1」乘给和，方向相反、方法相同。

\solve $x,y>0$：
\[
x+y=(x+y)\left(\frac1x+\frac9y\right)=1+9+\frac yx+\frac{9x}{y}\ge 10+2\sqrt{\frac yx\cdot\frac{9x}y}=10+6=16,
\]
当且仅当 $\dfrac yx=\dfrac{9x}y$，即 $y=3x$；代入条件 $\dfrac1x+\dfrac{9}{3x}=\dfrac4x=1$ 得 $x=4,\ y=12$。最小值 $16$。

\checkup $x=4,y=12$：条件 $\dfrac14+\dfrac34=1$ 成立，$x+y=16$。

\insight{条件即「1」时直接乘。展开后的固定流程：常数打底、互倒对给出 $2\sqrt{\cdot}$、取等条件回代条件解出具体位置，三步缺一不可。}
\end{solution}

\begin{solution}{题 17 \starred\starred\ · 解析（双换元）}
\think 两个分母 $x+3y$、$x-y$ 都是 $x,y$ 的组合，配凑无从下手，于是设它们为新变量。换元后条件成为「和 $\le 4$」，不是等式；但目标为正、且随分母增大而减小，方向恰好相合。分离出 $\left(\dfrac2m+\dfrac1n\right)(m+n)$ 这一 0 次齐次部分（拓展 C 的视角），即可处理。

\solve 设 $m=x+3y,\ n=x-y$。由 $x>y>0$ 得 $m>0,\ n>0$。反解：
\[
x=\frac{m+3n}{4},\qquad y=\frac{m-n}{4},\qquad x+y=\frac{m+n}{2}.
\]
条件 $x+y\le 2$ 即 $m+n\le 4$。于是
\[
\frac{2}{x+3y}+\frac{1}{x-y}=\frac2m+\frac1n=\frac{\left(\frac2m+\frac1n\right)(m+n)}{m+n}\ \ge\ \frac{\left(\frac2m+\frac1n\right)(m+n)}{4},
\]
（最后一步用了 $m+n\le 4$ 且分子为正。）而
\[
\left(\frac2m+\frac1n\right)(m+n)=2+1+\frac{2n}{m}+\frac{m}{n}\ge 3+2\sqrt2,
\]
故原式 $\ge\dfrac{3+2\sqrt2}{4}$。

等号要求两处同时成立：$\dfrac{2n}m=\dfrac mn$（即 $m=\sqrt2\,n$）且 $m+n=4$。解得 $n=4(\sqrt2-1)$，$m=8-4\sqrt2$，回代：
\[
x=\frac{m+3n}{4}=2\sqrt2-1,\qquad y=\frac{m-n}{4}=3-2\sqrt2.
\]
验证范围：$x\approx1.83>y\approx0.17>0$，$x+y=2\le 2$，均满足。故最小值为 $\dfrac{3+2\sqrt2}{4}$。

\checkup 代回原式：$x+3y=8-4\sqrt2$，$\dfrac{2}{8-4\sqrt2}=\dfrac{4+2\sqrt2}{8}$；$x-y=4\sqrt2-4$，$\dfrac{1}{4\sqrt2-4}=\dfrac{\sqrt2+1}{4}$。相加得 $\dfrac{6+4\sqrt2}{8}=\dfrac{3+2\sqrt2}{4}\approx1.457$，一致。

\insight{分母是变量组合时，设分母为新变量。「和 $\le$ 定值」与「随分母增大而减小」的目标相配时同样可以取等，但等号出现在边界上，两处取等条件必须同时兑现。与题 16 对照：那里的「1」是现成的，这里的 $\frac{m+n}{4}\le 1$ 是自己构造的，并且带方向。}
\end{solution}

\clearpage

% ===========================================================================
%  课后练习 · 简解与答案（练习 1–13）
% ===========================================================================
\section*{课后练习\quad 简解与答案}

\begin{solution}{练习 1 · 简解}
由 $0<a<b,\ a+b=1$ 知 $a<\dfrac12$；$ab<\dfrac14$（等号需 $a=b$，取不到），故 $2ab<\dfrac12$；$a^2+b^2=(a+b)^2-2ab=1-2ab>\dfrac12$。四个数中只有 $a^2+b^2$ 超过 $\dfrac12$。\whisper{答案：$a^2+b^2$。}
\end{solution}

\begin{solution}{练习 2 · 简解}
正向成立：$ab\le\left(\dfrac{a+b}2\right)^2\le\left(\dfrac42\right)^2=4$。反向不成立，反例 $a=8,\ b=\dfrac14$：$ab=2\le 4$ 而 $a+b=8\dfrac14>4$。\whisper{答案：能推出；不能推出（反例如上）。}
\end{solution}

\begin{solution}{练习 3 · 简解}
A：$x<0$ 时不成立（$x=-1$ 时左边为 $-2$）。B：$\sqrt{x^2+2}\ge\sqrt2$ 恒成立。C：$x^2+\dfrac1{x^2}\ge 2>\sqrt2$ 恒成立。D：$x>0$ 时 $-3x-\dfrac4x<0$（如 $x=1$ 时左边为 $-5$），不成立。\whisper{答案：AD。}
\end{solution}

\begin{solution}{练习 4 · 简解}
展开：$(x+y)\left(\dfrac1x+\dfrac4y\right)=1+4+\dfrac yx+\dfrac{4x}y\ge 5+2\sqrt4=9$，等号 $y=2x$。\whisper{答案：$9$。}
\end{solution}

\begin{solution}{练习 5 \starred\ · 简解}
令 $t=x-1\in(-5,0)$：$y=\dfrac{t^2+1}{2t}=\dfrac12\left(t+\dfrac1t\right)$。$t<0$ 时 $t+\dfrac1t\le -2$（对 $-t>0$ 用基本不等式后还原符号），故 $y\le -1$，等号 $t=-1$ 即 $x=0\in(-4,1)$。左端点 $-4$ 不影响结论。\whisper{答案：最大值 $-1$（$x=0$）。}
\end{solution}

\begin{solution}{练习 6 \starred\ · 简解}
$m\le(2a+b)\left(\dfrac2a+\dfrac1b\right)$ 恒成立，等价于 $m\le$ 右边最小值。展开：$4+1+\dfrac{2a}b+\dfrac{2b}a=5+2\left(\dfrac ab+\dfrac ba\right)\ge 5+4=9$，等号 $a=b$。\whisper{答案：$m$ 的最大值为 $9$。}
\end{solution}

\begin{solution}{练习 7 · 简解}
$1=\dfrac x3+\dfrac y4\ge 2\sqrt{\dfrac{xy}{12}}$，故 $\dfrac{xy}{12}\le\dfrac14$，$xy\le 3$，等号 $\dfrac x3=\dfrac y4=\dfrac12$，即 $x=\dfrac32,\ y=2$。\whisper{答案：$3$。}
\end{solution}

\begin{solution}{练习 8 \starred\ · 简解}
$(x+2)+(y+1)=4$。乘 $1=\dfrac{(x+2)+(y+1)}{4}$：
\[
\frac{4}{x+2}+\frac{1}{y+1}=\frac14\left(5+\frac{4(y+1)}{x+2}+\frac{x+2}{y+1}\right)\ge\frac14(5+4)=\frac94,
\]
等号 $x+2=2(y+1)$，配合 $x+y=1$ 得 $x=\dfrac23,\ y=\dfrac13$。\whisper{答案：$\dfrac94$。}
\end{solution}

\begin{solution}{练习 9 · 简解}
（书例 1.2-1）改写 $8x=4(2x-1)+4$：
\[
f(x)=4(2x-1)+\frac{1}{2x-1}+4\ge 2\sqrt{4(2x-1)\cdot\frac1{2x-1}}+4=8,
\]
等号 $4(2x-1)=\dfrac1{2x-1}$，即 $(2x-1)^2=\dfrac14$，$2x-1=\dfrac12>0$，$x=\dfrac34$。\whisper{答案：$8$（$x=\frac34$）。}
\end{solution}

\begin{solution}{练习 10 · 简解}
（书例 1.3-4）条件两边同除 $5xy$：$\dfrac{1}{5y}+\dfrac{3}{5x}=1$。乘 1：
\[
3x+4y=(3x+4y)\left(\frac{1}{5y}+\frac{3}{5x}\right)=\frac{3x}{5y}+\frac{12y}{5x}+\frac95+\frac45\ge\frac25\sqrt{36}+\frac{13}5=5,
\]
等号 $\dfrac{3x}{5y}=\dfrac{12y}{5x}$ 即 $x=2y$，回代条件得 $y=\dfrac12,\ x=1$。\whisper{答案：C（$5$）。}
\end{solution}

\begin{solution}{练习 11 \starred\ · 简解}
（书例 1.2-2）A 错：$x=y=1$ 满足条件而 $x+y=2$。D 错：$x=\dfrac{\sqrt3}3,\ y=-\dfrac{\sqrt3}3$ 满足条件而 $x^2+y^2=\dfrac23<1$。C 对：由 $xy\le\dfrac{x^2+y^2}2$（\starref{}，对一切实数成立），$1=x^2+y^2-xy\ge\dfrac{x^2+y^2}2$，故 $x^2+y^2\le 2$。B 对：配方 $1=\dfrac{(x+y)^2}4+\dfrac{3(x-y)^2}4\ge\dfrac{(x+y)^2}4$，故 $(x+y)^2\le 4$，$-2\le x+y\le 2$。\whisper{答案：BC。}
\end{solution}

\begin{solution}{练习 12 \starred\starred\ · 简解}
（书例 1.2-3）$1=(5x^2+y^2)\cdot y^2$，两边乘 4：
\[
4=(5x^2+y^2)\cdot 4y^2\le\left[\frac{(5x^2+y^2)+4y^2}{2}\right]^2=\frac{25}4(x^2+y^2)^2,
\]
故 $(x^2+y^2)^2\ge\dfrac{16}{25}$，$x^2+y^2\ge\dfrac45$。等号：$5x^2+y^2=4y^2=2$，即 $y^2=\dfrac12,\ x^2=\dfrac3{10}$，代回条件 $y^2(5x^2+y^2)=\dfrac12\times2=1$ 成立。\whisper{答案：$\dfrac45$。}
\end{solution}

\begin{solution}{练习 13 \starred\ · 简解}
（书例 1.2-7）由条件 $z=x^2-3xy+4y^2$：
\[
\frac z{xy}=\frac{x^2+4y^2}{xy}-3\ge\frac{2\sqrt{x^2\cdot4y^2}}{xy}-3=4-3=1,
\]
等号 $x=2y$。此时 $z=4y^2-6y^2+4y^2=2y^2$，于是
\[
x+2y-z=2y+2y-2y^2=4y-2y^2=2-2(y-1)^2\le 2,
\]
$y=1$ 时取到（$x=2,\ z=2$，均为正数，且 $x^2-3xy+4y^2=4-6+4=2=z$ 成立）。\whisper{答案：C（$2$）。}
\end{solution}

\clearpage


% ---- 方法一览（自总结版可注释）----
% ===========================================================================
%  方法一览（自总结版：把「要点」一栏留空，或在 main.tex 注释本文件）
% ===========================================================================
\section*{方法一览}

\begin{center}
\begin{tabularx}{\linewidth}{@{}l L l@{}}
\toprule
\multicolumn{1}{l}{\color{indigo}方法} &
\multicolumn{1}{l}{\color{indigo}要点} &
\multicolumn{1}{l}{\color{indigo}出处} \\
\midrule
作差法 & 作差、通分、判断符号（糖水不等式） & 0.1、题 1 \\
\starref{} $\to$ 基本不等式 & 平方非负 $\Rightarrow a^2+b^2\ge 2ab$ $\Rightarrow$ 用 $\sqrt a$ 替换 $a$ & 卡片 1 \\
几何解释 & 半弦不过半径；一张图呈现四个平均数 & 题 5、拓展 A \\
证明不等式 & 逐项放缩，同向相乘，取等条件须一致（会师） & 题 6、7 \\
求最值 & 一正、二定、三相等 & 题 8--13 \\
配凑 & 改写为「分母 + 分母的倒数 + 常数」；配系数凑定值 & 题 10、12，练习 9 \\
条件最值 & 乘「1」代换；和积互化；双换元 & 题 14--17 \\
齐次化 & 变形围绕次数进行；0 次齐次式只依赖比值 & 拓展 C \\
对称性 & 用于猜测取等点，不能代替证明 & 拓展 D \\
\bottomrule
\end{tabularx}
\end{center}

\begin{strand}
学生自总结版：把「要点」一栏留空，听完课自行填写。
\end{strand}


\end{document}
